\documentclass[11pt]{article}

\usepackage[margin=1in]{geometry}
\usepackage[T1]{fontenc}
\usepackage[utf8]{inputenc}
\usepackage{booktabs}
\usepackage{longtable}
\usepackage{array}
\usepackage{xcolor}
\usepackage{hyperref}
\usepackage{listings}
\usepackage{amsmath}
\usepackage{siunitx}

\hypersetup{
  hidelinks,
  pdftitle={Multi-Structure MORPHEUS Manual},
  pdfauthor={Merlino Development Team}
}

\definecolor{MerlinoSlate}{HTML}{4A5561}
\definecolor{MerlinoBg}{HTML}{F7F8F6}
\emergencystretch=2em

\lstdefinestyle{merlino}{
  basicstyle=\ttfamily\small,
  breaklines=true,
  columns=fullflexible,
  keepspaces=true,
  frame=single,
  rulecolor=\color{MerlinoSlate!45},
  backgroundcolor=\color{MerlinoBg},
  xleftmargin=0.5em,
  xrightmargin=0.5em
}

\newcommand{\program}{\texttt{MORPHEUS Multi-Structure}}
\newcommand{\merlino}{\texttt{Merlino}}
\newcommand{\morpheus}{\texttt{MORPHEUS}}
\newcommand{\gicforge}{\texttt{GICForge}}
\newcommand{\zeff}{Z_\mathrm{eff}}

\title{\textbf{Multi-Structure MORPHEUS Manual}\\
\large Shared class-correction refinement across related molecular structures}
\author{Merlino Development Notes}
\date{\today}

\begin{document}
\maketitle

\begin{abstract}
\program{} is the ensemble-refinement layer for related molecules, conformers
or homologous series.  Instead of refining every molecule independently, it
fits transferable corrections to classes of structural coordinates relative to
their own quantum-chemical reference geometries.  The workflow is designed for
systems where isotopic substitution is sparse for each molecule but chemical
class information is shared across the set.  This manual describes the model,
the TOML job format, class selectors, soft priors, hard constraints, synthon
\(\zeff\) atom typing, diagnostics and recommended validation workflows.
\end{abstract}

\tableofcontents
\newpage

\section{Purpose}

Single-molecule semiexperimental refinement uses isotopologues of one molecule
to determine a non-redundant structural model.  Multi-structure refinement uses
several related molecules to determine shared corrections to chemically
equivalent coordinate classes.  The absolute structural parameters of different
molecules are not forced to be equal.  What is shared is the correction to the
computational reference for a declared coordinate class.

This is useful for homologous series, conformer pairs and large isotope-poor
molecules.  For example, two steroid hormones need not have equal C--C bond
lengths, but a transferable bias in the computed C--C skeleton class may be
estimated from rotational data for both molecules at once.  The same idea can
separate C--C, C--O single, C=O, angular and torsional classes.  C--H
distances can be retained at the high-level reference geometry, for example
BDPCS3, when the available isotopic data mainly constrain the heavy-atom frame.

\section{Theoretical Background}

For molecule \(m\) and generated coordinate \(k\), the model is
\[
  q_m^\mathrm{SE}(k)
  =
  q_m^\mathrm{QC}(k)
  +
  \Delta_{c(m,k)} ,
\]
where \(q_m^\mathrm{QC}(k)\) is the quantum-chemical reference value,
\(q_m^\mathrm{SE}(k)\) is the semiexperimental value implied by the ensemble
fit, and \(c(m,k)\) is the class assigned to that coordinate.  The fitted
unknowns are the class corrections \(\Delta_c\), not the full set of Cartesian
coordinates of every molecule.

Each molecule keeps its own geometry, topology, isotopologues, vibrational
corrections, symmetry and active coordinate model.  The ensemble layer
linearizes each molecule around its reference structure and builds one global
weighted design matrix for the shared class corrections.  If \(r_m\) is the
single-molecule spectroscopic residual and \(J_m\) is the derivative of that
residual with respect to the molecule's active coordinates, the class
projection \(C_m\) maps global corrections into molecule-specific coordinate
changes:
\[
  \delta q_m = C_m \Delta .
\]
The global weighted least-squares problem is therefore assembled from blocks
\[
  r_m + J_m C_m \Delta .
\]
Soft priors add Gaussian observations of the form
\[
  \Delta_c = \Delta_c^{(0)} \pm \sigma_c ,
\]
whereas hard constraints remove or fix selected corrections.  This gives three
natural variants: no prior, soft prior and hard constraint.  Comparing these
variants is a core diagnostic because poorly supported broad classes may be
rank-consistent but numerically ill-conditioned.

\section{Ensemble Job Format}

An ensemble job uses schema:

\begin{lstlisting}[style=merlino]
schema = "merlino.semiexp.ensemble.v1"
title = "maleic, phthalic and succinic anhydrides shared class-correction fit"
\end{lstlisting}

\noindent
The stable top-level sections are:

\begin{longtable}{@{}p{0.22\linewidth}p{0.68\linewidth}@{}}
\toprule
Section & Meaning \\
\midrule
\endhead
\texttt{[fit]} & Numerical finite-difference step and rank threshold. \\
\texttt{[acceptance]} & Rank, conditioning, support and correlation policy. \\
\texttt{[[molecules]]} & List of ordinary single-molecule MORPHEUS/SE jobs. \\
\texttt{[[classes]]} & Transferable coordinate-class definitions, optional
priors and optional synthon selectors. \\
\bottomrule
\end{longtable}

\subsection{Fit and acceptance policy}

\begin{lstlisting}[style=merlino]
[fit]
step = 1.0e-4
rcond = 1.0e-10

[acceptance]
require_full_rank = true
max_condition_number = 1.0e8
min_residual_degrees_of_freedom = 1
min_molecule_support = 2
high_correlation_review_threshold = 0.98
high_correlation_reject_threshold = 0.9999
\end{lstlisting}

\noindent
The result status is \texttt{accepted}, \texttt{review} or \texttt{rejected}.
\texttt{accepted} means the model passes rank, conditioning, molecule-support
and correlation checks.  \texttt{review} means the model is numerically usable
but contains strongly correlated corrections.  \texttt{rejected} means the
class model should not be used without changing classes, priors or atom typing.

\subsection{Molecule list}

\begin{lstlisting}[style=merlino]
[[molecules]]
name = "maleic_anhydride"
job = "../maleic_anhydride/maleic_anhydride_ensemble.mse.toml"

[[molecules]]
name = "phthalic_anhydride"
job = "../phthalic_anhydride/phthalic_anhydride_ensemble.mse.toml"
\end{lstlisting}

\noindent
Each molecule points to an ordinary single-molecule job.  That job owns its
geometry, isotopologues, vibrational corrections, coordinate model and
single-molecule options.  The ensemble layer does not merge geometries; it
merges only the class-correction parameters.

\section{Class Selectors}

A class must refer to one coordinate type.  The implementation rejects class
definitions that mix stretches with bends, torsions or out-of-plane
coordinates.  Selectors can use coordinate kind, atom symbols, value windows,
substrings, primitive signatures or synthon \(\zeff\) values.

\subsection{Atom-symbol selectors}

Atom-symbol matching is chemically canonicalized by coordinate type:

\begin{itemize}
  \item stretches use an unordered atom pair;
  \item bends preserve the central atom and canonicalize the two terminal atoms;
  \item torsions use the unordered central-atom pair, i.e. the chemistry of the
        central bond;
  \item out-of-plane coordinates use the central atom.
\end{itemize}

\begin{lstlisting}[style=merlino]
[[classes]]
name = "CO_carbonyl"
kind = "stretch"
atoms = ["C", "O"]
value_max = 1.28

[[classes]]
name = "CO_single"
kind = "stretch"
atoms = ["C", "O"]
value_min = 1.28
\end{lstlisting}

\noindent
Value windows are evaluated on the reference primitive values and allow broad
symbol classes to be split into chemically distinct subclasses.

\subsection{Soft priors and hard constraints}

Soft priors are declared directly on classes:

\begin{lstlisting}[style=merlino]
[[classes]]
name = "CCO_bend"
kind = "bend"
atoms = ["C", "C", "O"]
prior_value = 0.0
prior_sigma = 1.0e-3
\end{lstlisting}

\noindent
They regularize weakly supported shared corrections without hiding the
spectroscopic residuals.  A hard-constraint comparison fixes selected
corrections, usually to zero, and removes them from the global variable set.
The recommended analysis for a new chemical family is to compare no-prior,
soft-prior and hard-constraint variants.

\section{Synthon \texorpdfstring{\(Z_\mathrm{eff}\)}{Zeff} Typing}

The preferred high-resolution atom typing is continuous synthon typing.  Each
atom is described by an effective atomic number \(\zeff\), and two atoms are
considered equivalent when their \(\zeff\) values differ by less than a
declared threshold.  For a primitive coordinate, the chemically relevant
ordered signature is used: unordered pair for stretches, central atom retained
for bends, central bond for torsions, and central atom for out-of-plane
coordinates.

\begin{lstlisting}[style=merlino]
[[classes]]
name = "Ccarb_Ocarbonyl"
kind = "stretch"
atoms = ["C", "O"]
synthon_zeff = [5.91, 7.865]
synthon_threshold = 0.035
\end{lstlisting}

\noindent
Discrete \texttt{synthon\_signatures} remain available for reproducibility
checks, but thresholded \(\zeff\) classes are more transferable across
conformers and homologous molecules.

\section{Command-Line Workflows}

\subsection{Production ensemble fit}

\begin{lstlisting}[style=merlino]
python -m merlino semiexp-ensemble \
  --job examples/semiexp/anhydrides_ensemble/anhydrides_ensemble.mse-ensemble.toml \
  --outdir working/semiexp/anhydrides_ensemble
\end{lstlisting}

\subsection{Prior comparison}

\begin{lstlisting}[style=merlino]
python -m merlino semiexp-ensemble-compare \
  --job examples/semiexp/anhydrides_ensemble/anhydrides_ensemble.mse-ensemble.toml \
  --outdir working/semiexp/anhydrides_prior_comparison \
  --soft-prior-sigma 1.0e-3
\end{lstlisting}

\noindent
This writes no-prior, soft-prior and hard-constraint variants and a summary
comparison.  It is the preferred diagnostic when deciding whether a class model
is chemically informative or merely numerically regularized.

\subsection{Synthon threshold scan}

\begin{lstlisting}[style=merlino]
python -m merlino semiexp-ensemble-synthon-scan \
  --job examples/semiexp/glycine_ensemble/glycine_conformers_synthon.mse-ensemble.toml \
  --outdir working/semiexp/glycine_synthon_scan \
  --thresholds 0.015,0.025,0.035,0.050
\end{lstlisting}

\subsection{Paper artifact regeneration}

\begin{lstlisting}[style=merlino]
python -m merlino semiexp-ensemble-paper \
  --job examples/semiexp/anhydrides_ensemble/anhydrides_ensemble.mse-ensemble.toml \
  --paper-dir doc/papers/ensemble_jpcl \
  --outdir working/semiexp/anhydrides_full_analysis
\end{lstlisting}

\noindent
This refreshes CSV/JSON diagnostics, prior-strength scan, leave-one-molecule-out
diagnostics and LaTeX fragments for the JPCL-style multi-structure paper.

\section{Output Files}

\begin{longtable}{@{}p{0.40\linewidth}p{0.50\linewidth}@{}}
\toprule
Output & Meaning \\
\midrule
\endhead
\texttt{ensemble\_class\_corrections.txt} & Human-readable correction report. \\
\texttt{ensemble\_class\_corrections.csv} & Class correction, uncertainty and
support summary. \\
\texttt{ensemble\_class\_report.csv} & Per-class diagnostics, priors and
acceptance information. \\
\texttt{ensemble\_molecule\_blocks.csv} & Per-molecule row counts, matched
coordinates and residual reduction. \\
\texttt{ensemble\_manifest.json} & Scientific manifest for the fitted class
model. \\
\texttt{ensemble\_covariance.csv} & Covariance matrix for class corrections. \\
\texttt{ensemble\_correlation.csv} & Correlation matrix for class corrections. \\
\texttt{run\_manifest.json} & Workflow manifest written by CLI entry points. \\
\bottomrule
\end{longtable}

The prior-comparison workflow also writes \texttt{ensemble\_prior\_comparison}
CSV/JSON files and a prior-strength scan.  The paper workflow writes the same
diagnostics plus generated LaTeX fragments.

\section{Diagnostics and Acceptance}

The minimum diagnostics for a production model are:

\begin{itemize}
  \item rank of the weighted design matrix;
  \item scaled condition number;
  \item residual degrees of freedom;
  \item molecule support for each class;
  \item class-correction standard errors;
  \item covariance and correlation matrices;
  \item weighted RMS before and after correction;
  \item leave-one-molecule-out transfer score when enough molecules are
        available.
\end{itemize}

\noindent
The leave-one-molecule-out score is the held-out WRMS after prediction divided
by the uncorrected held-out WRMS.  Values below one indicate transferability to
the held-out molecule; values near or above one indicate that the class model is
not predictive for that molecule.

\section{Recommended Regression Set}

Current scientific regression targets are:

\begin{itemize}
  \item maleic, phthalic and succinic anhydrides: no-prior, soft-prior and
        hard-constraint comparisons;
  \item glycine conformers I/II: continuous synthon \(\zeff\) atom typing and
        threshold scan;
  \item primitive signature canonicalization: torsions matched by central bond
        and out-of-plane coordinates by central atom;
  \item GUI job-writer round trip for fine selectors, soft priors and absolute
        paths;
  \item CLI manifest generation for scientific and workflow manifests.
\end{itemize}

\begin{lstlisting}[style=merlino]
pytest merlino_semiexp/tests/test_ensemble.py \
  gui/tests/test_ensemble_window_helpers.py
\end{lstlisting}

\section{Relation to MORPHEUS}

\program{} is an extension of \morpheus{}, not a replacement for the
single-molecule workflow.  Single-molecule MORPHEUS constructs and refines the
best model supported by one molecule's isotopic information.  Multi-structure
MORPHEUS asks a different question: which corrections to a quantum-chemical
reference are transferable across a chemically related set?  The two workflows
share coordinate generation, topology checks, primitive projections,
constraints and manifests, but the fitted variables are different.

\end{document}
